\documentclass[11pt]{article}
\usepackage{icmlko}
\begin{document}

\title{사냥터에는 유령이 살았다 --- 도메인별 가림이 잰 것의 $70\%$ 는 이웃이고, 손해의 $78\%$ 는 정보가 아니라 배선이었다}
\author{월드모델 실험실}
\date{노트 525 · 2026년 8월 4일}
\maketitle

\begin{abstract}
새 규약(``이 채택이 여는 갈래'')의 둘째 항목 --- \texttt{TRAINCAP} $\cdot$
\texttt{TRAINW} $\cdot$ \texttt{TRAINR} 이 들어간 뒤 노트 503 이 죽인
\textbf{도메인별 축 가림} 후보들이 되살아났나. 인구조사를 지금 챔피언에서
다시 돌리고, 이번에는 \emph{후보 자체}를 직접 쟀다.

\begin{center}
\begin{tabular}{lrrr}
\hline
계열 & 노트 503 & 지금(503 셈법) & 지금(덮음 $\ge 1\%$) \\
\hline
위키 & $+0.0012$ & $+0.0011$ & $+0.0002$ \\
태그 & $+0.0014$ & $+0.0017$ & $\mathbf{+0.0000}$ \\
gen & $+0.0020$ & $+0.0011$ & $+0.0002$ \\
grp & $+0.0022$ & $+0.0010$ & $+0.0010$ \\
\hline
합 & $+0.0068$ & $+0.0050$ & $\mathbf{+0.0015}$ \\
\hline
\end{tabular}
\end{center}

\textbf{사냥터의 $70\%$ 가 이웃 효과다.} 태그는 $100\%$ 가 그렇다 --- 태그를
가진 도메인은 하나도 안 벌고, 버는 도메인은 태그 덮음이 $0\%$ 다. 그리고
남은 후보들도 \textbf{인구조사가 부호를 못 맞힌다}(여섯 중 셋이 반대).
가림 손해를 갈라 보니 \textbf{$78\%$ 가 정보가 아니라 배선}이었고, 배선을
고쳐도 판은 안 샀다. \textbf{사냥터는 진짜 비어 있다.}
\end{abstract}

\begin{figure}[t]
\centering
\includegraphics[width=\textwidth]{figs/f1.pdf}
\caption{\textbf{가} --- 인구조사는 여섯 다 양수라 했다. \textbf{나} ---
$\diamond$ 가림에서 $\circ$ 섞기를 빼면 배선의 몫. \textbf{다} --- 요약.}
\end{figure}

\section{인구조사는 순서만 맞다}

노트 503 의 인구조사는 한 계열을 \emph{전 도메인}에서 빼고 도메인별 차를
읽는다. 그런데 후보는 \emph{그 도메인만} 가리는 것이다. 앞은 자기 몫과
이웃 몫의 합이고 뒤는 자기 몫뿐이다. 덮음 $\ge 1\%$ 인 여섯 후보를
직접 쟀다(짝 씨앗 12).

\begin{center}
\begin{tabular}{lrrr}
\hline
도메인 $\cdot$ 계열 & 인구조사 & \textbf{자기만} & 판 차 \\
\hline
세계애니 grp & $+0.0070$ & $\mathbf{+0.0050}$ & $+0.0007$ \\
게임 gen & $+0.0033$ & $+0.0004$ & $-0.0010$ \\
모바일 위키 & $+0.0019$ & $+0.0006$ & $-0.0008$ \\
모바일 grp & $+0.0016$ & $\mathbf{-0.0039}$ & $-0.0024$ \\
웹툰 grp & $+0.0009$ & $\mathbf{-0.0120}$ & $-0.0025$ \\
펀딩 grp & $+0.0002$ & $\mathbf{-0.0177}$ & $-0.0047$ \\
\hline
\end{tabular}
\end{center}

순위상관은 $+0.943$ 이라 \emph{사전 등록은 통과했다}. 그런데

\begin{quote}
\textbf{사전 등록이 틀린 것을 쟀다.} 순위상관은 \emph{부호}에 둔감하다.
여섯이 전부 양수로 적혔는데 실제로는 \textbf{셋이 크게 음수}이고, 펀딩
grp 은 $+0.0002$ 로 적힌 것이 실제로는 $-0.0177$ --- \textbf{$88$ 배 틀리고
부호도 반대}다.
\end{quote}

거기에 \textbf{더해지지도 않는다}: 자기만 $+$ 나머지만 이 인구조사와
평균 $-0.0073$ 어긋난다. 펀딩 grp 은 자기 $-0.0177$ $+$ 나머지 $+0.0002$
인데 전부 빼면 $+0.0002$ 다. \textbf{grp 을 다 없애면 펀딩이 멀쩡하고,
펀딩에서만 없애면 펀딩이 무너진다.}

\section{손해는 정보가 아니라 배선이다}

가림은 두 가지를 한꺼번에 한다 --- ① 그 축의 정보를 없애고 ② 그 행을 남이
학습한 갈래의 \textbf{결측 쪽}으로 보낸다. \texttt{DOMPERM} 을 새로 달아
①만 하는 팔(도메인 안에서 그 축의 행 짝만 섞기 --- 값의 분포도 결측 무늬도
그대로)을 만들고 둘을 뺐다.

\begin{center}
\begin{tabular}{lrrr}
\hline
 & 가림 & 섞기 & \textbf{배선의 몫} \\
\hline
펀딩 grp & $-0.0177$ & $-0.0039$ & $\mathbf{-0.0137}$ $(78\%)$ \\
웹툰 grp & $-0.0120$ & $\mathbf{+0.0024}$ & $-0.0144$ \\
세계애니 grp & $+0.0050$ & $+0.0048$ & $-0.0002$ \\
애니 trend & $+0.0241$ & $+0.0203$ & $\mathbf{+0.0038}$ \\
\hline
\end{tabular}
\end{center}

웹툰 grp 은 \textbf{부호가 뒤집힌다} --- 정보만 없애면 오히려 벌고,
가림이 해로워 보인 것은 전부 배선 탓이었다.

\subsection*{그런데 배선을 고쳐도 판은 안 산다}

섞기 팔의 판 차는 웹툰 $-0.0002$($6/12$) $\cdot$ 세계애니 $-0.0002$($5/12$)
로 \textbf{$0$ 을 문다}. 도메인은 벌지만 이웃이 그만큼 잃는다.

\begin{quote}
\textbf{사냥터는 진짜 비어 있다.} 겉보기 $+0.0050$ 중 $70\%$ 는 이웃 효과,
남은 것의 절반은 부호가 반대, 그리고 도구(가림)의 결함을 고쳐도 판은
안 산다. \textbf{도메인별 축 가림 계열은 닫는다.}
\end{quote}

\section{배선의 부호는 축마다 다르다}

펀딩 grp 은 배선이 $-0.0137$ 로 해롭고 애니 trend 는 $+0.0038$ 로
\emph{이롭다}. 학습 행에서 그 축의 전역 결측을 보면

\begin{center}
\begin{tabular}{lrr}
\hline
계열 & 전역 덮음 & 결측 \\
\hline
grp & $68\%$ & $\mathbf{32\%}$ \\
gen & $66\%$ & $34\%$ \\
태그 & $54\%$ & $46\%$ \\
trend\_ & $40\%$ & $\mathbf{60\%}$ \\
wiki\_ & $7\%$ & $93\%$ \\
\hline
\end{tabular}
\end{center}

\begin{quote}
\textbf{읽기(관측으로 둔다)} --- 결측 행이 적은 축(grp $32\%$)을 도메인별로
가리면 그 행이 \emph{안 배워진} 결측 갈래로 떨어져 무너진다. 결측 행이 많은
축(trend\_ $60\%$)은 결측 갈래가 잘 배워져 있어 오히려 안전하고, 도메인을
갈라 주는 일까지 한다.
\end{quote}

$n=2$ 라 법칙이라 하지 않는다. \textbf{반증 예측을 함께 적는다}(노트 503 이
남긴 조항): 태그($46\%$)와 gen($34\%$)은 grp 쪽처럼, wiki\_($93\%$)는
trend\_ 쪽처럼 행동해야 한다.

\section*{조항}

\begin{quote}
\textbf{도메인별 축 가림 후보는 둘을 함께 본다} --- ① 그 도메인의 그 축
\textbf{덮음}($0\%$ 면 이웃 효과라 후보가 아니다, 노트 503) ② 그 축의
\textbf{전역 결측 비율}(낮으면 결측 갈래가 안 배워져 있어 가림이 배선으로
무너진다, 여기).
\end{quote}

그리고 도구에 관한 것 하나 더 --- \textbf{훑기가 후보의 값을 재고 있는지
먼저 확인한다.} 노트 503 의 인구조사는 ``전부 빼고 도메인 차''라 후보와
\emph{다른 것}을 쟀고, 순위만 맞아서 $22$ 노트 동안 그 사실이 안 드러났다.

\section*{남은 것}

\begin{itemize}
\item \textbf{\texttt{DOMPERM} 은 남겨 둔다}(기본 빈 사전 = 챔피언 불변).
섞기가 가림보다 나은 자리가 어디 또 있는지 아직 안 훑었다.
\item \textbf{애니 trend 는 섞기로 바꾸면 안 된다} --- 가림 $+0.0241$ 대
섞기 $+0.0203$, 판으로도 $-0.0015$. 지금 채택이 맞다.
\item \textbf{판 정의} --- 이 노트의 모든 판 수치는 유보 행수 가중이다.
도메인 단순 평균으로 바꾸면 어떻게 되는지 아직 모른다(노트 526 에서 잰다).
\end{itemize}

\end{document}
